\section{Hiện thực hệ thống}

\subsection{Nền tảng sử dụng}

\subsubsection{Python 3.11 và gRPC}

\textbf{Python 3.11} là ngôn ngữ chính cho AI Pipeline. Lý do chọn Python:
hệ sinh thái ML/AI phong phú nhất (\texttt{sentence-transformers},
\texttt{pdfplumber}, \texttt{openai}); cộng đồng lớn; asyncio native hỗ trợ
I/O-bound workload hiệu quả.

\textbf{gRPC + Protobuf:} Giao thức giao tiếp giữa các service. Ưu điểm so
với REST/JSON: binary encoding nhanh hơn $\sim$3--5x, schema strongly typed
qua \texttt{.proto}, hỗ trợ streaming response.

\subsubsection{MinIO Object Storage}

MinIO là hệ thống object storage mã nguồn mở tương thích S3 API. Trong đề
tài, MinIO lưu trữ file tài liệu gốc (PDF, DOCX, PPTX). Bucket chính:
\texttt{documents}. Mỗi file có key: \texttt{\{owner\_id\}/\{doc\_id\}/original.\{ext\}}.

\subsubsection{Qdrant Vector Database}

Qdrant chạy trong Docker container, persistent volume tại
\texttt{./data/qdrant}. Collection \texttt{document\_chunks}: vector 768
chiều, metric cosine, payload gồm \texttt{document\_id}, \texttt{chunk\_id},
\texttt{heading}, \texttt{language}.

\subsubsection{React + TypeScript --- Frontend}

Frontend được xây dựng bằng \textbf{React 18 + TypeScript + Vite}. Thư viện
UI: \textbf{Tailwind CSS} + \textbf{shadcn/ui} (headless components). Giao
tiếp với Backend qua \textbf{gRPC-Web} (\texttt{@grpc/grpc-js} + Envoy proxy).

\subsubsection{Docker và Docker Compose}

\begin{lstlisting}[language=bash, caption={Cấu trúc \texttt{docker-compose.yml} (rút gọn)}]
services:
  ai-pipeline:    # Python gRPC server
  frontend:       # React dev server
  qdrant:         # Vector database
  minio:          # Object storage
  postgres:       # Relational database
  envoy:          # gRPC-Web proxy

networks:
  internal:       # Service-to-service communication
\end{lstlisting}

\subsection{Hiện thực AI Pipeline}

\subsubsection{Module Document Parsing}

Ba parser được triển khai, mỗi cái implement \texttt{ParserPort}:

\begin{lstlisting}[language=Python, caption={Interface ParserPort}]
from abc import ABC, abstractmethod
from src.domain.entities.document import Document, ParsedDocument

class ParserPort(ABC):
    @abstractmethod
    def parse(self, document: Document,
              file_bytes: bytes) -> ParsedDocument:
        """Trích xuất text và cấu trúc heading từ file."""
        ...
\end{lstlisting}

\begin{lstlisting}[language=Python, caption={PDFParser sử dụng pdfplumber}]
class PDFParser(ParserPort):
    def parse(self, document: Document,
              file_bytes: bytes) -> ParsedDocument:
        sections = []
        with pdfplumber.open(BytesIO(file_bytes)) as pdf:
            for page in pdf.pages:
                text = page.extract_text() or ""
                # Phát hiện heading theo font size/bold
                headings = self._detect_headings(page)
                sections.extend(
                    self._build_sections(text, headings, page.page_number)
                )
        return ParsedDocument(document_id=document.id,
                              sections=sections)
\end{lstlisting}

\subsubsection{Module Text Chunking}

\begin{lstlisting}[language=Python, caption={HeadingChunker --- xử lý section dài}]
class HeadingChunker(ChunkerPort):
    def chunk(self, parsed_doc: ParsedDocument) -> list[Chunk]:
        chunks = []
        for section in parsed_doc.sections:
            token_count = self._count_tokens(section.text)
            if token_count <= self.max_chunk_size:
                chunks.append(self._to_chunk(section))
            else:
                # Sliding window với overlap
                sub_chunks = self._sliding_window(
                    section, self.max_chunk_size, self.overlap
                )
                chunks.extend(sub_chunks)
        return chunks

    def _sliding_window(self, section: Section,
                        size: int, overlap: int) -> list[Chunk]:
        sentences = sent_tokenize(section.text)
        windows, current, current_tokens = [], [], 0
        for sent in sentences:
            t = self._count_tokens(sent)
            if current_tokens + t > size and current:
                windows.append(self._join(section, current))
                # Giữ lại phần overlap
                current = current[-(overlap // 20):]
                current_tokens = sum(self._count_tokens(s)
                                     for s in current)
            current.append(sent)
            current_tokens += t
        if current:
            windows.append(self._join(section, current))
        return windows
\end{lstlisting}

\subsubsection{Module sinh nội dung AI}

\begin{lstlisting}[language=Python, caption={OpenAIAdapter --- sinh Summary và Quiz}]
class OpenAIAdapter(LLMPort):
    def __init__(self, api_key: str, model: str = "gpt-4o-mini"):
        self.client = AsyncOpenAI(api_key=api_key)
        self.model = model

    async def generate_summary(self, text: str,
                                language: str = "vi") -> list[str]:
        prompt = SUMMARY_PROMPT.format(
            text=text, language=language, n_points=5
        )
        response = await self.client.chat.completions.create(
            model=self.model,
            response_format={"type": "json_object"},
            messages=[{"role": "user", "content": prompt}],
            temperature=0.3,
        )
        data = json.loads(response.choices[0].message.content)
        return data["summary_points"]

    async def generate_quiz(self, text: str,
                             language: str = "vi") -> list[dict]:
        prompt = QUIZ_PROMPT.format(
            text=text, language=language, n_questions=5
        )
        response = await self.client.chat.completions.create(
            model=self.model,
            response_format={"type": "json_object"},
            messages=[{"role": "user", "content": prompt}],
            temperature=0.5,
        )
        data = json.loads(response.choices[0].message.content)
        return data["questions"]
\end{lstlisting}

\subsubsection{Use Case ProcessDocument}

\begin{lstlisting}[language=Python, caption={ProcessDocumentUseCase --- điều phối pipeline}]
class ProcessDocumentUseCase:
    async def execute(self, command: ProcessDocumentCommand):
        doc = await self.doc_repo.get(command.document_id)
        # 1. Parse
        await self._update_status(doc.id, DocumentStatus.PARSING)
        file_bytes = await self.storage.download(doc.file_key)
        parser = self.parser_factory.get(doc.file_type)
        parsed_doc = parser.parse(doc, file_bytes)
        # 2. Chunk
        await self._update_status(doc.id, DocumentStatus.CHUNKING)
        chunks = self.chunker.chunk(parsed_doc)
        await self.chunk_repo.save_all(chunks)
        # 3. Embed + Generate (chạy theo batch để tránh rate limit)
        await self._update_status(doc.id, DocumentStatus.GENERATING)
        for batch in batched(chunks, size=5):
            await asyncio.gather(*[
                self._process_chunk(chunk) for chunk in batch
            ])
        await self._update_status(doc.id, DocumentStatus.DONE)

    async def _process_chunk(self, chunk: Chunk):
        # Embedding
        vector = await self.embedder.embed(chunk.content)
        await self.vector_db.upsert(str(chunk.id), vector,
                                    chunk.metadata)
        # Summary + Quiz song song
        summary, quiz = await asyncio.gather(
            self.llm.generate_summary(chunk.content),
            self.llm.generate_quiz(chunk.content),
        )
        await self.card_repo.save(Card(chunk_id=chunk.id,
                                       summary=summary,
                                       quiz_questions=quiz))
\end{lstlisting}

\subsection{Hiện thực Web Application}

\subsubsection{Giao diện Giảng viên}

\begin{figure}[H]
    \centering
    % \includegraphics[width=0.9\textwidth]{figures/screen_dashboard.png}
    \caption{Màn hình Dashboard Giảng viên.}
    \label{fig:screen-dashboard}
\end{figure}

\begin{figure}[H]
    \centering
    % \includegraphics[width=0.9\textwidth]{figures/screen_upload.png}
    \caption{Màn hình Upload tài liệu.}
    \label{fig:screen-upload}
\end{figure}

\begin{figure}[H]
    \centering
    % \includegraphics[width=0.9\textwidth]{figures/screen_progress.png}
    \caption{Màn hình theo dõi tiến trình xử lý AI (4 bước).}
    \label{fig:screen-progress}
\end{figure}

\begin{figure}[H]
    \centering
    % \includegraphics[width=0.9\textwidth]{figures/screen_review.png}
    \caption{Màn hình Review nội dung AI (layout 2 cột).}
    \label{fig:screen-review}
\end{figure}

\subsubsection{Giao diện Học viên}

\begin{figure}[H]
    \centering
    % \includegraphics[width=0.85\textwidth]{figures/screen_courses.png}
    \caption{Màn hình danh sách khóa học của Học viên.}
    \label{fig:screen-courses}
\end{figure}

\begin{figure}[H]
    \centering
    % \includegraphics[width=0.85\textwidth]{figures/screen_card.png}
    \caption{Màn hình xem bài học dạng Card.}
    \label{fig:screen-card}
\end{figure}

\begin{figure}[H]
    \centering
    % \includegraphics[width=0.85\textwidth]{figures/screen_quiz.png}
    \caption{Màn hình làm Quiz (hiển thị phản hồi tức thì sau khi chọn đáp án).}
    \label{fig:screen-quiz}
\end{figure}

\subsection{Tổng kết Chương 5}

Chương 5 đã trình bày chi tiết việc hiện thực hệ thống: lý giải lựa chọn
từng nền tảng công nghệ, mã nguồn cốt lõi của AI Pipeline (Parser, Chunker,
LLM Adapter, Use Case) và giao diện Web Application cho cả hai nhóm người dùng.
Toàn bộ mã nguồn được tổ chức theo kiến trúc DDD, đảm bảo tính modular và
dễ kiểm thử --- cơ sở để tiến hành kiểm thử toàn diện trong Chương 6.
